% =============================================================================
% frontmatter.tex  --  title page, scope/authorities, how-to-read + rule key, ToC
% =============================================================================

% ---- title page -----------------------------------------------------------
\begin{titlingpage}
  \centering
  \vspace*{\stretch{1.2}}
  {\Huge\scshape The Construction of\\[3pt] the FTD Mathematics\par}
  \vspace{14pt}
  {\large\itshape From a Discrete Ontology to Its Provable Boundary\par}
  \vspace{20pt}
  {\rule{0.46\textwidth}{0.4pt}\par}
  \vspace{20pt}
  {\small\scshape a monograph on foundational ternary dynamics\par}
  \vspace{\stretch{1}}
  \begin{minipage}{0.82\textwidth}
    \centering\itshape\small
    Derive everything the discrete ontology forces --- and rigorously
    establish the boundary of what it cannot force.
  \end{minipage}\par
  \vspace{\stretch{1.4}}
  % NOTE (O-1): author line deliberately omitted. The source monograph carries
  % no author; the project's grade-A paper leaves \author{} empty for the owner
  % to fill at submission time. Do NOT fabricate a name.
  {\small 2 June 2026\par}
  \vspace*{\stretch{0.4}}
\end{titlingpage}

% ---- scope ----------------------------------------------------------------
\section*{Scope}
\noindent This monograph builds FTD's mathematics from the ground up, in the
order a mathematician would build the number systems, for an interdisciplinary
and philosophy-of-physics readership. It is a construction and an honest
accounting at once: it introduces no new mathematics, and it states every result
at exactly the grade the result has earned --- theorem, derivation, conjecture,
or open problem --- and at no grade higher. That discipline is not a hedge; it is
the point. It lets the strong results stand as strong and the boundary stand as a
genuine finding.

% ---- how to read ----------------------------------------------------------
\section*{How to read this monograph}
\noindent This is a \emph{construction story}. It builds FTD's mathematics in
the order a mathematician would build the number systems --- each object from
prior ones, each step independently checkable --- and it asks one question at
every step: \textbf{does the discrete ontology force this, or not?} (where
``the discrete ontology'' includes the one $\mathbb{Z}[i]$ modelling choice of
\secref{sec:seed}). The answer comes with an epistemic tag, and the tag is the
point. A reader from philosophy of physics is invited to scrutinize the tags
first and the prose second; the tags are where the project's honesty lives.

\medskip
\noindent The monograph is organized around the project's single governing
goal, stated as two clauses:
\begin{quote}\itshape
  Derive everything the discrete ontology forces --- and rigorously establish
  the boundary of what it cannot force.
\end{quote}
\noindent Both clauses are deliverables. Showing that the discrete ontology
\emph{does not} fix a quantity is as much a result of this program as showing
that it does; a mapped boundary is not a failure but a finding.

\medskip
\noindent\textbf{How a tag looks.} Throughout, an epistemic tag is set as a
small-caps badge whose \emph{rule} --- never its colour --- carries its tier,
so the page reads identically in grayscale. There are four treatments:
\rulekeylegend
\noindent The same rule signatures frame the displayed ``marquee'' statements
and label the construction map of \cref{fig:dag}. The full contract --- every
tag, its meaning, and what a reviewer should expect of it --- is the table of
\secref{sec:contract}.

\cleardoublepage
\tableofcontents*
